\documentclass{article}
\usepackage{amsmath, amssymb, amsthm}
\usepackage{hyperref}
\usepackage{geometry}
\geometry{margin=1in}

\title{Erd\H{o}s Problem \#138}
\date{Last updated: 2025-08-31}
\author{Source: erdosproblems.com}

\begin{document}
\maketitle

\noindent\textbf{Status:} OPEN \quad \textbf{Prize: \$500}

\noindent\textbf{Tags:} additive combinatorics

\medskip

\noindent\textbf{Problem Statement:}

Let the van der Waerden number $W(k)$ be such that whenever $N\geq W(k)$ and $\{1,\ldots,N\}$ is $2$-coloured there must exist a monochromatic $k$-term arithmetic progression. Improve the bounds for $W(k)$ - for example, prove that $W(k)^{1/k}\to \infty$.




When $p$ is prime Berlekamp \cite{Be68} has proved $W(p+1)\geq p2^p$. Gowers \cite{Go01} has proved\[W(k) \leq 2^{2^{2^{2^{2^{k+9}}}}}.\]The best general lower bound is $W(k)\gg 2^k$, due to Kozik and Shabanov \cite{KoSh16}. In \cite{Er81} Erd\H{o}s further asks whether $W(k+1)/W(k)\to \infty$, or $W(k+1)-W(k)\to \infty$. DeepMind has proved (see the comments) that\[W(k+1)\geq W(k)+k,\]answering the second question. In \cite{Er80} Erd\H{o}s asks whether $W(k)/2^k\to \infty$, and offers \$500 for a proof or disproof of $W(k)^{1/k}\to \infty$. The analogous question for $r\geq 3$ colours was resolved by Fox and Hunter \cite{FoHu26}, who proved\[W_3(k)^{1/k}\geq C^{\log_*k}\]for some constant $C&#62;1$.




References

[Be68] Berlekamp, E. R., A construction for partitions which avoid long arithmetic progressions . Canad. Math. Bull. (1968), 409-414.

[Er80] Erd\H{o}s, Paul, A survey of problems in combinatorial number theory . Ann. Discrete Math. (1980), 89-115.

[Er81] Erd\H{o}s, P., On the combinatorial problems which I would most like to see solved . Combinatorica (1981), 25-42.

[FoHu26] J. Fox and Z. Hunter, Three-color van der Waerden numbers grow super-exponentially . arXiv:2606.02541 (2026).

[Go01] Gowers, W. T., A new proof of Szemer\'{e}di's theorem . Geom. Funct. Anal. (2001), 465-588.

[KoSh16] Kozik, Jakub and Shabanov, Dmitry, Improved algorithms for colorings of simple hypergraphs and applications . J. Combin. Theory Ser. B (2016), 312--332.

\noindent\textbf{Related OEIS sequences:} A005346


\bigskip
\noindent\small{Source: \url{https://www.erdosproblems.com/138}}
\end{document}
